\section{Conclusion}

This report started to investigate the transformative integration of Artificial Intelligence and Large Language Models into software engineering through a Systematic Literature Review that mapped current benefits, consequences, and stakeholder opinions. From the information collected in the SLR, it could be shown that while AI adoption is accelerating across the industry, organizations currently lack structured methodologies to integrate these tools responsibly without compromising software quality, security, or accountability. To address this gap, this research is proposing a new framework based on Agile and Scrum principles, designed to formalize Human-AI collaboration and guide the governance of AI-enabled development workflows.

The primary limitations of this study relate to the scope of its evaluation strategy. As the framework is assessed through a demonstration within a specific consultancy context and relies on qualitative expert interviews, the generalization of results to broader organizational settings may be constrained. Additionally, because the evaluation focuses on the initial stages of adoption, long-term impacts such as sustained organizational change and lasting performance improvements cannot yet be fully measured.

Future work will focus on the comprehensive development and refinement of the proposed framework, followed by its practical demonstration and "ex post" evaluation within the next-generation consultancy environment. This process will assess the artifact’s validity, practical suitability, and perceived value through the outlined interviews and success criteria. The research will culminate in the final dissertation and the submission of scientific papers, aiming to provide actionable standards for the responsible era of AI-augmented software development.

